%=============================================================================!
% 2.A  ANSWERS TO THE CHECKS                                                  !
%=============================================================================!
\unnumbered{Answers to the checks}
\label{sec:ne-answers}

\answerto{sec:ne-force}Zero. The puck moves in a straight line
at constant speed, so its velocity is constant, and by the first law the
net force on it vanishes.

\answerto{sec:ne-frames}Taking forwards as positive, the
passenger's velocity in the train is $v' = \SI{-1}{\metre\per\second}$,
so adding $u$ to both sides of \cref{eq:ne-velocity-frames}, $v' = v -
u$, the platform sees $v = v' + u = -1 + 5 = \SI{4}{\metre\per\second}$
forwards. A roundabout is not an inertial
frame of reference: each point on it moves in a circle and so accelerates
towards the centre (\cref{sec:mo-circle}), and a ball rolled across it
appears, to someone riding it, to curve with no sideways force.

\answerto{sec:ne-second}$a = F/m = 6/2 =
\SI{3}{\metre\per\second\squared}$. Doubling the mass to
\SI{4}{\kilogram} halves it, to \SI{1.5}{\metre\per\second\squared}.

\answerto{sec:ne-third}The two forces are equal, but the
acceleration each produces is the force divided by the mass of the body
it acts on. The apple's mass is about \SI{0.1}{\kilogram} and the Earth's
about \SI{6e24}{\kilogram}, so the apple accelerates at $g$ and the Earth
by an amount far too small to detect.

\answerto{sec:ne-gravity}In a vacuum they land together,
because the only force on each is its weight, and $\vb{a} = \vb{g}$
whatever the mass. In air the hammer lands first: the drag on the feather
is comparable to its tiny weight, while the drag on the hammer is
negligible beside its weight.

\answerto{sec:ne-equations}No: both start at rest, so their
velocities agree, but their positions differ by \SI{1}{\metre}, so their
states differ. Their motions differ too, the lower stone staying
\SI{1}{\metre} below the upper at every instant. Matching velocities
alone does not fix the motion.

\answerto{sec:ne-drag}It slows down. Above the terminal
velocity the drag is larger than the weight, so the net force points
upwards. By \cref{eq:ne-drag-v} the gap $v_0 - v_\infty$ is positive and
shrinks exponentially, so the velocity falls towards $v_\infty$ from
above, steeply at first and then ever more gently, never crossing the
dotted line of \cref{fig:ne-drag}(a).

\answerto{sec:ne-spring}By \cref{eq:ne-spring-omega} the
period is proportional to $\sqrt{m}$, so four times the mass doubles it,
to \SI{2}{\second}. The period does not depend on the amplitude, so
pulling twice as far leaves it at \SI{1}{\second}.

\answerto{sec:ne-bodies}The skater and ball start at rest, so
their total momentum is zero. On smooth ice no horizontal external force
acts, the weight and the push of the ice being vertical and balancing,
so by \cref{eq:ne-momentum} the horizontal momentum stays zero: $60\,v +
0.5\times12 = 0$, giving $v = -6/60 = \SI{-0.1}{\metre\per\second}$. The skater slides backwards at
\SI{0.1}{\metre\per\second}.

\answerto{sec:ne-atoms}The lithium atom, which is lighter. Under
one force the acceleration is inversely proportional to the mass, so it
accelerates $15.999/6.94 = 2.31$ times faster than the oxygen atom.
